\documentclass[10pt]{article}

\usepackage{amsmath}
\usepackage{amssymb}
\usepackage{amsthm}
\usepackage{ mathrsfs }
\usepackage{ mathtools}
\usepackage{enumerate}
\usepackage{bbm}
\usepackage{lipsum}
\usepackage{fancyhdr}
\usepackage{tikz-cd} 
\usepackage{adjustbox} 
\usetikzlibrary{arrows}
\newcommand{\midarrow}{\tikz \draw[-triangle 90] (0,0) -- +(.1,0);}
\tikzset{commutative diagrams/.cd,
mysymbol/.style={start anchor=center,end anchor=center,draw=none}
}
\newcommand\comsymb[2][\circlearrowleft]{%
  \arrow[mysymbol]{#2}[description]{#1}}

\newtheorem{theorem}{Theorem} [section] 
\newtheorem{proposition}{Proposition}[section] 
\newtheorem{definition}{Definition}[section] 
\newtheorem{lemma}{Lemma}[section] 
\newtheorem{notation}{Notation}[section] 
\newtheorem{remark}{Remark}[section] 
\newtheorem{corollary}{Corollary} [section] 
\newtheorem{terminology}{Terminology}[section] 
\newtheorem{fact}{Fact}[section] 
\newtheorem{example}{Example}[section] 
\newtheorem{claim}{Claim}
\newenvironment{claimproof}[1]{\par\noindent\underline{Proof:}\space#1}{\hfill $\blacksquare$}

\DeclareMathOperator{\triv}{triv}
\DeclareMathOperator{\ind}{ind}
\DeclareMathOperator{\straw}{s.t.}
\DeclareMathOperator{\ev}{ev}
\DeclareMathOperator{\pr}{pr}
\DeclareMathOperator{\Cl}{Cl}
\DeclareMathOperator{\Aut}{Aut}
\DeclareMathOperator{\Map}{Map}
\DeclareMathOperator{\Stab}{Stab}
\DeclareMathOperator{\Ind}{Ind}
\DeclareMathOperator{\GL}{GL}
\DeclareMathOperator{\SL}{SL}
\DeclareMathOperator{\SO}{SO}
\DeclareMathOperator{\Sp}{Sp}
\DeclareMathOperator{\esssup}{esssup}
\DeclareMathOperator{\abs}{abs}
\DeclareMathOperator{\rel}{rel}
\DeclareMathOperator{\diam}{diam}
\DeclareMathOperator{\rank}{rank}
\DeclareMathOperator{\Hom}{Hom}
\DeclareMathOperator{\Homk}{Hom_{k-alg}}
\DeclareMathOperator{\Ker}{Ker}
\DeclareMathOperator{\Image}{Im}
\DeclareMathOperator{\Dom}{Dom}
\DeclareMathOperator{\grad}{grad}
\DeclareMathOperator{\divergence}{div}
\DeclareMathOperator{\rk}{rk}
\DeclareMathOperator{\Span}{Span}
\DeclareMathOperator{\MaxSpec}{MaxSpec}
\DeclareMathOperator{\interior}{int}
\DeclareMathOperator{\supp}{supp}
\DeclareMathOperator{\id}{id}
\DeclareMathOperator{\sgn}{sgn}
\DeclareMathOperator{\Mat}{Mat}
\DeclareMathOperator{\PD}{PD}
\DeclareMathOperator{\ext}{ext}
\DeclareMathOperator{\vol}{vol}
\DeclareMathOperator{\inte}{int}
\DeclareMathOperator{\diverg}{div}
\DeclareMathOperator{\curl}{curl}
\DeclareMathOperator{\image}{im}
\DeclareMathOperator{\dR}{dR}
\DeclareMathOperator{\Ob}{Ob}
\DeclareMathOperator{\Mnfd}{Mnfd}
\DeclareMathOperator{\OpEmb}{OpEmb}
%\newcommand*{\name}[\num_arguments][default values]{{\color{#1}\Large #2}}
\newcommand{\defeq}{\vcentcolon=}
\newcommand{\norm}[1]{\Vert #1 \Vert}
\newcommand{\opNorm}[2]{\norm{#1}_{#2\to#2}}
\newcommand{\normL}[3]{\norm{#1}_{L^{#2}(#3)}}
\newcommand{\N}[0]{\mathbb{N}}
\newcommand{\m}[0]{\mathfrak{m}}
\newcommand{\R}[0]{\mathbb{R}}
\newcommand{\Z}[0]{\mathbb{Z}}
\newcommand{\C}[0]{\mathbb{C}}
\newcommand{\Q}[0]{\mathbb{Q}}
\newcommand{\F}[0]{\mathbb{F}}
\newcommand{\G}[0]{\mathbb{G}}
\newcommand{\A}[0]{\mathbb{A}}
\newcommand{\Para}[0]{\mathbb{P}}
\newcommand{\M}[0]{\mathcal{M}}
\newcommand{\maniN}[0]{\mathcal{N}}
\newcommand{\cinf}[0]{\mathcal{C}^\infty}
\newcommand{\torus}[0]{\mathbb{T}}
\newcommand{\sep}[0]{\:|\:}
\newcommand{\pd}[2]{{\frac {\partial #1} {\partial #2}}}
\newcommand{\compinc}[0]{\subset \subset}
\newcommand{\sH}[0]{\mathscr{H}}
\newcommand{\st}[0]{\:\straw\:}
\newcommand{\fib}[1]{%
  \mathbin{\mathop{\times}\limits_{#1}}%
}
\newcommand{\tens}[1]{%
  \mathbin{\mathop{\otimes}\displaylimits_{#1}}%
}

\title{Algebraic Number Theory}
\author{Lectures given by\\Dr. Alexander Horawa\\Typed by\\So Murata}
\date{SoSe 26, University of Bonn}


\begin{document}
\maketitle

%4/13

\section{Motivation of Algebraic Number Theory}

The goal of the field is to solve the rational polynomial equations. 
\begin{example}
    Finding all Pythagorean triples $(x,y,z)$. i.e. a triplet $(x,y,z)$ satisfying,
    \begin{equation*}
    x^2+y^2=z^2.
    \end{equation*}
    The obstacle here is that the equation is stated in both products and sums.
    To overcome this problem, we consider,
    \begin{equation*}
    (x+iy)(x-iy) = z^2.
    \end{equation*}
    The equation holds in the ring $\Z[i]$. In this form, we will be able to consider the prime factors of $z$. This is justified by the following proposition.
\end{example}
\begin{proposition}
    $\Z[i]$ is a UFD.
\end{proposition}
\begin{proof}
    Exercise.
\end{proof}
\begin{corollary}
    $\Z[i]^\times = \{\pm1,\pm i\}$.
\end{corollary}
\begin{proof}
    Exercise.
\end{proof}
\begin{lemma}
    \label{lemma_prime_divisor_z}
    Suppose $\pi\in\Z[i]$ be an irreducible such that for a primitive Pythagorean triplet $(x,y,z)$ we have $\pi\mid z$. Then $\pi$ divides exactly one of $x+iy$ and $x-iy$ so does $\pi^2$.
\end{lemma}
\begin{proof}
    If $\pi$ divides $(x+iy)$ and $(x-iy)$ simultaneously, we get a contradiction as the triplet is primitive.
\end{proof}
\begin{corollary}
    \label{corollary_factor_z_square}
    In the above setting, we have,
    \begin{equation*}
    x+ iy = u\alpha^2
    \end{equation*}
    where $\alpha\in\Z[i]$ and $u\in\Z[i]^\times$.
\end{corollary}
\begin{proof}
    By Lemma \ref{lemma_prime_divisor_z}, the statement clearly follows. 
\end{proof}
\begin{theorem}
    All Pythagorean triples $(x,y,z)$ are of the following form.
    \begin{equation*}
    (x,y,z) = (\pm(m^2-n^2),\pm2mn,\pm(m^2+n^2))
    \end{equation*}
\end{theorem}
\begin{proof}
    By Corollary \ref{corollary_factor_z_square}, suppose $\alpha = m+in$ for some $m,n\in\Z$. Without loss of generality, we assume $u=1$. We have,
    \begin{equation*}
        \alpha^2 = m^2-n^2+2imn,\overline{\alpha}^2 = m^2-n^2-2imn.
    \end{equation*}
    The statement then follows.
\end{proof}
Another important example would be the following celebrated theorem.
\begin{theorem}[Fermat-Wiles]
    For all $n\in\N$, if $n\geq3$ then there does not exist a triplet $(x,y,z)\in (\Z\backslash\{0\})^3$ such that 
    \begin{equation*}
        x^n+y^n = z^n.
    \end{equation*}
\end{theorem}
The case where $n=3,4$ are elementary. Also the statement is equivalent to say that it holds for all prime numbers greater or equal to $3$. We will go through Kummer's attempt to prove Fermat's Last Theorem.
\begin{notation}
A $n$-th primitive root will be denoted by $\omega_n$. 
\end{notation}
\begin{remark}
    Two important rings are
    \begin{equation*}
        \Z[\omega_p]\subseteq\Q(\omega_p).
    \end{equation*}
    As we have seen it in Galois theory, $\Q(\omega_p)$ is a splitting field of $x^p-1$ over $\Q$.
\end{remark}
Over $\Z[\omega_p]$, we have the following factorization,
\begin{equation*}
    x^p+y^p = \prod_{i=0}^{p-1}(x+\omega_p^iy).
\end{equation*}
\begin{remark}
    If $\Z[\omega_p]$ is a UFD, then there are no solutions to FLT. 
\end{remark}
However, $\Z[\omega_p]$ is rarely a UFD. We will see this explicitly in Homework 1.
\begin{theorem}
There are finitely many primes $p$ such that $\Z[\omega_p]$ is a UFD.
\end{theorem}
Kummer's main idea was that even though $\Z[\omega_p]$ fails to be a UFD, we can factor ideals of $\Z[\omega_p]$ uniquely into prime ideals. The ideal version of the factorization would be the following.
\begin{equation*}
    (x+y)(x+\omega_py)\cdots(x+\omega^{p-1}y) = (z)^p,
\end{equation*}
then there is an ideal $I\subseteq\Z[\omega_p]$ such that 
\begin{equation*}
    (x+\omega_py)=I^p.
\end{equation*}
If $I$ happens to be principle, (i.e. $I=(\alpha)$) then we have,
\begin{equation*}
    x+\omega_py = u\cdot\alpha^p
\end{equation*}
where $u\in\Z[\omega_p]^\times$. 
\par However, the above attempt also runs into a problem. That is $\Z[\omega_p]$ may not be a PID. In general we consider,
\begin{enumerate}
    \item A field $K$ which is a fintie extension of $\Q$ which we will call the Number Field.
    \item $\mathcal{O}_K$ which is the set of integral elements of $K$ which we will call the Ring of Integers.
\end{enumerate}
\begin{example}
    For $K=\Q(\omega_p)$, we have $\mathcal{O}_K = \Z[\omega_p]$.
\end{example}
The classs group measures how far ideals are from being principal.
\begin{definition}
\begin{equation*}
    \Cl(K)\defeq \{\text{ideals of $\mathcal{O}_K$}\}/\{\text{principal ideals of $\mathcal{O}_K$}\}.
\end{equation*}
\end{definition}
\begin{remark}
    The class group is not necessarily abelian. 
\end{remark}
The following statements are obvious.
\begin{enumerate}
    \item An ideal $I$ of $\mathcal{O}_K$ is principale if and only if $[I]\in\Cl(K)$ is trivial.
    \item $\mathcal{O}_K$ is a PID if and only if $\Cl(K) = \{1\}$.
\end{enumerate}
\begin{theorem}
    $\Cl(K)$ is finite.
\end{theorem}
\begin{definition}
    For a number field $K$, we define the class number of $K$ as
    \begin{equation*}
        h(K)\defeq\#\Cl(K).
    \end{equation*}
\end{definition}
\begin{remark}
    By Lagrange's theorem, we have the following identity. For any ideal $I$ of $\mathcal{O}_K$, $I^{h(K)}$ is principal.
\end{remark}
\end{document}